\section{Discussion}\label{sec:discussion}

\subsection{What the analysis already implies for operators}
Four consequences of \cref{sec:analysis} hold whatever the measurements show, and they shape how the measured trade-offs should be read.
First, padding cannot hide participation or the number of rounds, because it changes neither message counts nor timing (\cref{prop:pad}). An operator who only pads protects the size of individual replies and nothing else.
Second, hiding participation has an energy floor at every site, in every session, including sites that hold nothing relevant (\cref{prop:cover}). Rounds with every site active make all links look the same within a session, but the session-level shape still depends on which sites are called. If participation is not sensitive in a deployment, cover sites can be switched off and the budget spent on size and timing. If it is sensitive, the network-wide cost grows with the number of sites, which matters most for many small edge sites.
Third, the DP release is the only primitive with a formal guarantee, and it is expensive: with the queue-length release its padding grows with the square of the number of ticks at a fixed privacy budget (\cref{eq:dpcost}), every site is active at every tick, and the guarantee needs a long-lived framed connection, a session-start signal to all sites and a fixed schedule (\cref{sec:deploy}). Releasing noisy arrivals makes padding grow only linearly in the number of ticks but leaves the backlog unbounded (\cref{rem:exponent}).
Fourth, primitives do not compose for free: padding before the DP release inflates the bytes that $\Delta$ must cover (\cref{prop:order}), so $\Delta$ has to be measured after padding.

\subsection{What the completed measurements support}
The observed cost ordering provides engineering guidance for this workload, not a privacy ranking. Padding adds modest byte cost in the tested configuration; rounds increase latency and cover traffic; fixed-tick DP release remains much more expensive in bytes, even with the lower biases chosen through functional calibration. The minimal TLS baseline shows that all-site controls account for much of the apparent unshaped baseline cost. Reporting only shaped/unshaped ratios without the control protocol would obscure this contribution.

The 168 successful EC2 sessions establish that the chosen workflow can execute through both relay and NetShaper paths under the recorded conditions. The two local models' 48 passing checks show compatibility with typed plans and structured summaries after a documented decimal-presentation repair. Neither observation establishes general agent accuracy, empirical metadata secrecy, a 10\,Gbps throughput result or measured energy savings. Resource budgets and correctness gates are useful deployment inputs, but a calibrated privacy--energy frontier requires additional measurements.

\subsection{Limitations and future work}
This is a design-and-analysis paper with bounded engineering validation. The energy coefficients remain symbolic; the four virtual machines expose no physical joule measurements. There is no edge-device study, no CO$_2$ calculation and no empirical validation of the proposed joule-budgeted controller. Bytes and active duration may inform a future energy calibration but cannot stand in for one. We therefore make no claim of privacy gain per joule or superiority at equal energy.

The formal model has separately observed site links, whereas the cloud validation has 20 logical partitions behind one physical data endpoint and a multiplexed tunnel. Its captures do not support per-site participation inference. The retrospective question-disjoint multiplexed-link diagnostic in \cref{sec:retrospective} uses a small acquired TRAIN subset; it does not replace a confirmatory attacker test, and the reserved final-test bank remains unopened. An attacker study needs a matching observer layout, retraining under each shaping configuration and frozen selection before final-test scoring. The DP theorem is conditional on its framing, cap, schedule, session-start and composition assumptions; a functioning transport alone does not certify the wire-level guarantee.

Only 12 selected question strata and two repetitions are used in each cloud mode, and mode phases are not interleaved. Time, burst capacity and execution order may confound differences. The public-data projection and hand-authored tasks do not represent all institutional workloads. The model checks use two local models and structured numerical JSON, not a hosted model, autonomous replanning or general natural-language answer quality. The cloud task driver is deterministic. Planning and final-summary traffic are outside the measured observer segment.

Packet-window agreement is sensitive to endpoint timestamp boundaries, as shown by the two ACKs in the rounds configuration. Session-window metrics exclude setup and inter-session cover, which are reported separately. The new loopback comparison instead counts the full TLS lifecycle, and shows why periodic-cover frequency must be varied. Its fixed application-frame savings remain schedule-dependent, and packet jitter is not a privacy certificate. Large-scale concurrency, distinct physical sites, instrumented power calibration and final-test attacker evaluation remain future work. The public-envelope and risk-allocation candidates still require atomic enqueue enforcement, distributed TLS lifecycle/capture integration and independent per-site wire acquisition; the synthetic and single-host results must not be presented as multi-host deployment gains. Chains of agents, compromised sites, active attacks and traffic to a hosted LLM API are outside the threat model.

\section{Conclusion}\label{sec:conclusion}
Green Shaping combines tool-layer padding, cover rounds and fixed-tick DP release with an explicit model of per-node energy budgets. Its analysis identifies limits on size protection, a participation-hiding energy floor, release privacy conditions, padding growth and finite-horizon completion limits. Public-contract bounds support a restricted resource-admission candidate, and a fixed workflow envelope exposes a conditional resource--latency trade-off. The proposed attacker- and energy-budgeted controller is a design contribution whose empirical optimum has not been established.

The completed implementation study reports functional TRAIN/SELECT calibration, 48 successful repaired local-model checks, 168 reference-correct four-instance cloud sessions, matched capture accounting and independently verified recovery. Retrospective question-disjoint classification and fixed-arrival mechanism stress tests add bounded evidence about metadata predictability and the reliability cost of finite horizons; independent-seed simulations expose failures missed by small pilot screens. The measured bytes and latency expose framing overhead and substantial DP release costs rather than demonstrating privacy or energy superiority. Preserving that distinction makes the artifact a reproducible starting point for future measurements on independently observed sites with calibrated energy instrumentation.
